\documentclass[11pt]{article}
\usepackage[margin=1in]{geometry}
\usepackage{amsmath,amssymb,amsthm,hyperref}
\usepackage{mathrsfs}
\newtheorem{theorem}{Theorem}
\newtheorem{lemma}[theorem]{Lemma}
\newcommand{\dd}{\mathrm d'\mathrm d''}
\newcommand{\MA}{\operatorname{MA}}
\newcommand{\R}{\mathbb R}
\title{Bounded Green identities and the full bounded dual of radial measures}
\author{Research working paper, continuation 388}
\date{9 October 2026}
\begin{document}
\maketitle
\begin{abstract}
We give a local energy-completion proof of Green symmetry for continuous
differences of plurisubharmonic functions and a positive closed current.
It respects the smallest-sheaf definition of continuous psh functions,
including its successive uniform-closure stages. On a compact analytic
surface this implies first variation, concavity, and the tangent inequality
on the entire ordinary bounded continuous potential class. Subject to the
explicit, still unreviewed analytic branch-calibration premises below,
the tangent inequality upgrades a radial restricted dual calculation to
the full bounded-class dual. For the specified discrete branch measures
that dual is finite exactly under a cubic-tail condition. This is a new
working proof, awaiting independent review. It does not solve the entire
all-direction unbounded theory or the original P16.7.
\end{abstract}

\section{Framework and the exact dependency boundary}
Let $X$ be a good, Hausdorff, pure $n$-dimensional analytic space without
boundary, with the ordinary smooth forms, smooth cutoffs and partitions
of unity, currents, positivity, and Stokes calculus of CLD. A positive
closed current $T$ below has degree $(n-1,n-1)$, so that $\dd(fT)$ has
top degree. Pairings involving its smooth coefficients are Radon
measures. It suffices that positivity holds on strongly positive test
forms: the gradient tests used below are of rank one, and their sums.
We do not identify the whole middle-degree positive cone with the
strong cone or assume a positive definite ambient order unit.

Write $\mathscr P$ for the smallest sheaf of continuous functions
containing smooth psh functions and closed under local uniform limits,
and $\mathscr P^{\pm}$ for its vector-space span as a sheaf. These are
the ordinary CLD continuous psh and dpsh sheaves. We use the already
defined bounded product $\dd(fT)$: for $f\in\mathscr P$ it is positive
closed, and for $f\in\mathscr P^{\pm}$ it is a signed Radon measure.
The product is defined by continuous multiplication and differentiation
of currents; its existence does not presume the Green identity proved
here. The source statement is CLD-v2 Lemma 8.5.5, with its full immediate
proof; the continuous multiplier statement is Lemma 8.5.1.

For the global application take $n=2$, $X$ compact, and a smooth closed
positive form $\theta$. Let
\[
 \mathcal C_\theta=\{u\in C(X):\rho+u\in\mathscr P
       \text{ locally whenever }\dd\rho=\theta\}.
\]
Local smooth potentials $\rho$ and the ordinary bounded mixed calculus
are explicit hypotheses. Then $u$ is locally dpsh and
$\theta_u=\theta+\dd u$ is positive closed. We use the multilinear,
symmetric bounded products, continuous coefficient multiplication, and
local smooth-factor rules already constructed in whole continuation 383.
Their source foundation is disclosed; this paper does not newly prove
the quotient geometry, calibration, or every recursive source theorem.
Green symmetry and the full-class tangent inequality are proved below,
and are not additional assumptions.

\section{Local energy completion}
For smooth real $(1,0)$-forms $a,b$ of compact support use the symmetric
pairing induced by
\[
             \int a\wedge Jb\wedge T,
\]
where $J(\mathrm d'f)=\mathrm d''f$ is the usual symmetry convention.
Positivity and symmetry of $T$ make this a positive semidefinite
bilinear form. Quotient its null space and complete. On an open patch
we use the corresponding local Hilbert seminorms with smooth nonnegative
compact cutoffs $\chi$:
\[
 Q_{T,\chi}(f,g)=\int\chi\,\mathrm d'f\wedge\mathrm d''g\wedge T.
\]
Cauchy--Schwarz follows by applying positivity to $a+tb$ for every
real $t$. Local Hilbert limits of gradients are understood consistently
on smaller patches. Multiplication by a continuous function bounded
there extends to this completion, by uniform approximation by smooth
functions and the bound $\|fa\|\leq\|f\|_\infty\|a\|$.

Call a continuous function energy-approachable on a patch if it has
local smooth approximants converging uniformly and whose gradients
converge strongly in these local Hilbert seminorms. This condition defines
a vector-space sheaf $\mathscr H_T$. In particular it is closed under
multiplication by smooth cutoffs. To patch finitely many local
approximations on a relatively compact set, take a smooth partition
$\chi_i$ and approximate by $\sum_i\chi_i f_{i,k}$. Its derivative is
$\sum_i\chi_i\mathrm d'f_{i,k}+\sum_i f_{i,k}\mathrm d'\chi_i$;
the first terms converge strongly and the second uniformly in each
local Hilbert seminorm. This proves the asserted sheaf property.
Products of locally bounded members are also in $\mathscr H_T$:
approximate both factors and use the same argument for the Leibniz rule.

The gradient of an energy-approachable function is unique. Indeed,
if smooth $f_k\to0$ uniformly and their gradients converge strongly,
then for any compactly supported smooth $(1,0)$-form $a$, closedness of
$T$ and Stokes give
\[
 \int\mathrm d'f_k\wedge Ja\wedge T
       =-\int f_k\,\mathrm d'(Ja)\wedge T\longrightarrow0.
\]
The last integral tends to zero because the smooth-factor pairing with
$T$ is a finite signed measure on the fixed compact support.
The limit gradient is orthogonal to every smooth test form, a dense
class in the Hilbert completion, so is zero. Cutoffs give the same
conclusion locally. We can therefore write $\nabla_T f$ and
$Q_{T,\chi}(f,g)$ for members of $\mathscr H_T$ without a choice of
approximants.

\begin{lemma}\label{lem:test}
Suppose $f\in\mathscr H_T$ and $L_Tf=\dd(fT)$ is a signed Radon
measure. For compactly supported $g\in\mathscr H_T$,
\[
                 \int g\,L_Tf=-Q_T(f,g).
\]
All the expressions are local to a relatively compact neighborhood of
the support of $g$. In particular, if $w\in\mathscr H_T$ has signed
Radon $L_Tw$ and $\chi\geq0$ is smooth with compact support, then
\begin{equation}\label{eq:local-energy}
 Q_{T,\chi}(w,w)=\frac12\int w^2\dd\chi\wedge T
                         -\int\chi w\,L_Tw.
\end{equation}
\end{lemma}
\begin{proof}
For smooth compact $g$ and smooth approximants $f_k$, distributional
differentiation, closedness of $T$, and Stokes give
\[
 \langle L_Tf,g\rangle=\int f\dd g\wedge T
    =\lim_k\int f_k\dd g\wedge T
    =-\lim_k Q_T(f_k,g)=-Q_T(f,g).
\]
Use compactly supported smooth approximants $g_k$ for $g$, obtained by
multiplying local approximants by a cutoff equal to one near its
support. The left side converges by uniform convergence and finite
variation of $L_Tf$; the right by strong gradient convergence. This
proves the first formula. It does not require bounds on the Laplacians
of the smooth approximants $g_k$.

Apply it to $f=w$, $g=\chi w$. The product rule in the Hilbert
completion gives
\[
 -\int\chi w\,L_Tw
     =Q_{T,\chi}(w,w)+\int w\,\mathrm d'w\wedge\mathrm d''\chi\wedge T.
\]
The last term equals $-\tfrac12\int w^2\dd\chi\wedge T$, by the
same smooth approximation, $\mathrm d'(w^2)=2w\mathrm d'w$, and
closedness of $T$. This is \eqref{eq:local-energy}.
\end{proof}

\begin{lemma}\label{lem:psh}
Every continuous psh function is locally energy-approachable. Hence
every continuous dpsh function is a section of $\mathscr H_T$.
\end{lemma}
\begin{proof}
Consider the sheaf of continuous members of $\mathscr H_T$ for which
$L_Tu\geq0$. It contains smooth psh functions. We verify closure under
uniform limits; minimality of $\mathscr P$ then proves the assertion,
including all successive uniform-closure stages.

Let $u_j$ belong to this sheaf and tend locally uniformly to $u$.
On a relatively compact neighborhood of the support of $\chi$, bound
all $|u_j|$ by $B$. Choose another smooth compact cutoff $\eta\geq0$
equal to one near that support. Positivity gives
\[
 \int\chi(L_Tu_j+L_Tu_k)
   \leq\|\chi\|_\infty\int\eta(L_Tu_j+L_Tu_k)
   \leq2B\|\chi\|_\infty\|\dd\eta\wedge T\|_{\rm TV}.
\]
Here $\int\eta L_Tu_j=\int u_j\dd\eta\wedge T$ follows directly
from the definition of the current. The bound is independent of $j,k$.
Set $w=u_j-u_k$, $\varepsilon=\|w\|_\infty$ on the fixed neighborhood.
By \eqref{eq:local-energy} and positivity of the two Laplacians,
\[
 0\leq Q_{T,\chi}(w,w)
  \leq\frac{\varepsilon^2}{2}\|\dd\chi\wedge T\|_{\rm TV}
       +2B\varepsilon\|\chi\|_\infty\|\dd\eta\wedge T\|_{\rm TV}.
\]
Thus their gradients are strongly Cauchy locally. On any smaller patch,
choose smooth approximants to $u_j$ with uniform error and gradient
error at most $1/j$ for a cutoff equal to one on that patch. A diagonal
choice gives smooth approximants to $u$, uniformly and strongly in the
local energy seminorm. This proves $u\in\mathscr H_T$. The currents
$L_Tu_j=\dd(u_jT)$ converge distributionally to $L_Tu$ and are
positive, so $L_Tu\geq0$. This verifies uniform closure of the sheaf.
Finally its vector-space span contains every dpsh function.
\end{proof}

\begin{theorem}\label{thm:green}
If $X$ is compact without boundary and $f,g$ are continuous dpsh
functions, then for every positive closed $T$ as above,
\[
 \int f\dd(gT)=\int g\dd(fT)=-Q_T(f,g),\qquad
                  \int f\dd(fT)\leq0.
\]
\end{theorem}
\begin{proof}
Lemma~\ref{lem:psh}, finite patching on the compact space, and
Lemma~\ref{lem:test} apply. The bounded dpsh calculus makes the two
Laplacians signed Radon measures. The Hilbert pairing is symmetric and
nonnegative on the diagonal. There is no boundary term. No global
smooth psh regularization, unbounded recognition theorem, or skeletal
quotient of the function space is used.
\end{proof}

\section{First variation and concavity on the whole bounded class}
Assume now $n=2$ and the exact global bounded premises of Section 1.
Use the ordinary energy, with its usual normalization,
\[
 E(u)=\frac13\int_Xu(\theta^2+\theta\wedge\theta_u+\theta_u^2).
\]
All products with the continuous coefficient $u$ are actual bounded
current products. Constants shift this energy by their value times
$\int_X\theta^2$; in the branch application the latter is two.

\begin{theorem}\label{thm:tangent}
For $u,v\in\mathcal C_\theta$, put $w=v-u$, $u_t=u+tw$. Then
\[
 \frac{\mathrm d}{\mathrm dt}E(u_t)=\int_Xw\,\theta_{u_t}^2,
 \qquad
 \frac{\mathrm d^2}{\mathrm dt^2}E(u_t)
     =2\int_Xw\dd(w\theta_{u_t})
     =-2Q_{\theta_{u_t}}(w,w)\leq0.
\]
Consequently, for every pair in the entire class,
\begin{equation}\label{eq:tangent}
             E(v)\leq E(u)+\int_X(v-u)\MA(u).
\end{equation}
If all bounded operators have common raw mass $M$, then
$|E(v)-E(u)|\leq M\|v-u\|_\infty$, and $E(v)\geq E(u)$
whenever $v\geq u$.
\end{theorem}
\begin{proof}
Multilinearity makes $E(u_t)$ an ordinary cubic polynomial, so there
is no issue differentiating a nonsmooth potential in the parameter.
Write $A=\theta_{u_t}$ and note that $\dd u_t=A-\theta$.
Differentiating the energy gives
\[
 \frac13\left[\int w(\theta^2+\theta A+A^2)
       +\int u_t\dd(w\theta)+2\int u_t\dd(wA)\right].
\]
Apply Theorem~\ref{thm:green} with $T=\theta$ and $T=A$ to the
last two terms. They become respectively
$\int w(A-\theta)\theta$ and $2\int w(A-\theta)A$.
The resulting sum is $3\int wA^2$. Differentiate again using
multilinearity to obtain $2\int w\dd(wA)$, and use the same Green
theorem with $f=g=w$. The segment remains in $\mathcal C_\theta$,
so $A$ is positive closed for every $0\leq t\leq1$. Concavity of
the polynomial and its derivative at zero prove \eqref{eq:tangent}.
Integrating the first variation, positivity and common mass give
the stated Lipschitz bound and monotonicity.
\end{proof}

\begin{theorem}\label{thm:bounded-comparison}
Under the retained locality and common-mass premises, every
$u,v\in\mathcal C_\theta$ satisfies the full bounded comparison
inequality
\[
             \MA(v)\{u<v\}\leq\MA(u)\{u<v\}.
\]
\end{theorem}
\begin{proof}
The maximum of two continuous psh functions is continuous psh,
as in CLD-v2 Corollary 8.2.3. For the smallest continuous sheaf
this can also be checked directly. For smooth psh $a,b$, the smooth
psh functions $k^{-1}\log(e^{ka}+e^{kb})$ converge uniformly to
$\max(a,b)$, with error at most $k^{-1}\log2$.
Fixing smooth $b$, the functions $a\in\mathscr P$ for which
$\max(a,b)\in\mathscr P$ form a sheaf containing the smooth psh
functions and closed under uniform limits; hence they include all
of $\mathscr P$. Now fix such an $a$ and repeat this sheaf argument
in $b$. This proves the assertion at every uniform-closure stage.
Therefore
$w_\varepsilon=\max(u,v-\varepsilon)$ belongs to
$\mathcal C_\theta$. Ordinary locality gives $\MA(w_\varepsilon)
=\MA(u)$ on the open set $\{u>v-\varepsilon\}$ and
$\MA(w_\varepsilon)=\MA(v)$ on $\{u<v-\varepsilon\}$.
All three measures have the same total mass $M$. Positivity implies
\[
 \MA(v)\{u<v-\varepsilon\}
 \leq M-\MA(u)\{u>v-\varepsilon\}
 =\MA(u)\{u\leq v-\varepsilon\}.
\]
Choose $\varepsilon_j\downarrow0$ so that
$\MA(u)\{u=v-\varepsilon_j\}=0$. Such a sequence exists:
disjoint level sets of the continuous $u-v$ can have positive mass
for only countably many levels of a finite measure. The preceding
right side then equals $\MA(u)\{u<v-\varepsilon_j\}$.
Continuity from below of the two finite measures proves the assertion.
No product on the contact set was presumed.
\end{proof}

\section{A full bounded dual for the specified radial measures}
We retain, as separately unreviewed geometric and calibration inputs,
the entire branch construction of whole 387 and the stated bounded
premises of 383--385. Here is the exact amount required. The space is
the actual compact analytic surface, raw mass two; $o$ is its specified
origin Gauss point, $y$ the rigid infinite-depth point, and
$\xi_t$ its actual continuous depth section. Fix $1\leq N_1<N_2<\cdots$
with $N_i\to\infty$. For any finitely supported $c_i\geq0$ with
$\sum c_i\leq1$, the actual bounded potential
\[
 u_m=u_0-g_m(T),\quad
 h_m(t)=\sum_{i\leq m}c_i\mathbf1_{[0,N_i)}(t),\quad
 g_m(t)=\int_0^t h_m(s)\,\mathrm ds
\]
belongs to the entire ordinary $\mathcal C_\theta$ above, and its
actual bounded operator and energy satisfy
\begin{equation}\label{eq:calibration}
 \MA(u_m)=(2-h_m(0)^2)\delta_o+\xi_*[-\mathrm d(h_m^2)],
 \qquad E(u_m)=-\frac16-\frac13\int h_m^3.
\end{equation}
On $o$ and every $\xi_t$ the reference value is $u_0=-1/6$,
and $g_m(T(o))=0$. The two-atom outside identification and the
calibrated analytic product proof are premises here, not newly verified
by the energy argument. They are not inverse definitions of an operator.

Take a possibly infinite sequence $c_i\geq0$ with sum $c\leq1$,
and put $h=\sum_i c_i\mathbf1_{[0,N_i)}$ and
\[
 Z=\int_0^\infty h^3,\qquad
 \mu_h=(2-c^2)\delta_o+\xi_*[-\mathrm d(h^2)].
\]
These are actual analytic Radon measures: the second term is a finite
positive Stieltjes measure of mass $c^2$. We do not need to presume
full unbounded intrinsic recognition of the limiting potential to
compute their bounded-class dual
\[
 E^*_{\mathcal C_\theta}(\mu)
          =\sup_{v\in\mathcal C_\theta}\left(E(v)-\int_Xv\,\mu\right).
\]

\begin{theorem}\label{thm:dual}
Under exactly the retained calibration premises,
\[
 E^*_{\mathcal C_\theta}(\mu_h)=\frac16+\frac23Z
\]
in $\R\cup\{+\infty\}$. Thus this is the full bounded-class dual,
not just a supremum over the radial subfamily, and it is finite if and
only if $Z<\infty$.
\end{theorem}
\begin{proof}
Write $Z_m=\int h_m^3\uparrow Z$ and $\mu_m=\MA(u_m)$.
The continuous coefficients are bounded, so Stieltjes layer cake and
Tonelli, including the atoms, give
\[
 \int_Xu_m\,\mu_m=-\frac13-Z_m,
 \qquad E(u_m)-\int_Xu_m\,\mu_m=\frac16+\frac23Z_m.
\]
Indeed $\int g_m\,[-\mathrm d(h_m^2)]=\int h_m^3$; all measures
have total mass two and the $o$ term of $g_m$ is zero.

The measures $\mu_m$ converge weakly to $\mu_h$. Their tails past
a fixed depth $R$ are at most $h(R)^2$, which tends to zero since
$h(t)\to0$. On bounded depth intervals use the Stieltjes convergence
of the monotone profiles; continuity of $t\mapsto\xi_t$ transfers
this to analytic continuous tests. The outside coefficient tends to
$2-c^2$. This proof gives no missing atom at $y$.

If $Z<\infty$, fix an arbitrary $v\in\mathcal C_\theta$.
Theorem~\ref{thm:tangent}, applied to the actual bounded $u_m$, yields
\[
 E(v)-\int v\,\mu_m\leq\frac16+\frac23Z_m.
\]
Pass to the limit using continuity of the bounded $v$ and weak
convergence. This holds for every member of the entire bounded class,
so proves the full upper bound.

For the reverse bound, use the bounded $u_m$ themselves as tests
against the fixed $\mu_h$. Tonelli gives
\[
 \int_Xu_m\,\mu_h=-\frac13-\int_0^\infty h_mh^2,
 \qquad
 E(u_m)-\int_Xu_m\,\mu_h
     =\frac16-\frac13Z_m+\int_0^\infty h_mh^2.
\]
When $Z$ is finite the last integral increases to $Z$, proving
the required lower bound. When $Z=+\infty$, it is still finite for
each $m$ since $h_m$ has compact support, and $h\geq h_m$ gives
\[
 E(u_m)-\int_Xu_m\,\mu_h\geq\frac16+\frac23Z_m\longrightarrow+\infty.
\]
This proves the infinite case without an expression of the form
infinity minus infinity.
\end{proof}

Equivalently let $p_i\geq0$, $P=\sum_i p_i\leq1$, and
$\mu=(2-P)\delta_o+\sum_i p_i\delta_{\xi_{N_i}}$.
Set $q_i=\sum_{j\geq i}p_j$ and
$c_i=\sqrt{q_i}-\sqrt{q_{i+1}}$. The coefficients telescope to
$\sqrt P\leq1$, and $-\mathrm d(h^2)$ has exactly the prescribed
atoms. Therefore
\[
 E^*_{\mathcal C_\theta}(\mu)
   =\frac16+\frac23\int_0^\infty
          \left(\sum_{N_i>t}p_i\right)^{3/2}\mathrm dt.
\]
For a finite collection truncate the notation in the evident way.
The endpoint convention does not affect the integral. If this integral
is finite, the generated canonical radial solution and its genuine
non-pluripolar product are the separately proved construction of 387.
This theorem now proves finite full bounded dual for its measure,
under the stated bounded and calibration inputs. It makes no assertion
that $P\leq1$ is necessary for solutions outside this particular family.

\section{Source ledger, review status, and remaining scope}
The new energy-completion, sheaf-closure, Green, full bounded variation,
and limiting dual arguments are given in full in this paper and remain
UNREVIEWED. No conclusion is assumed from an active native campaign;
no active mathematical output was read or fed back. All original
counts stay at 16 solved, 26 open, and 97 previously accepted scoped
papers. Original P16.7 remains open.

For the continuous psh definition ROOT freshly read the complete
CLD-v2 Section 8.1 through Proposition 8.1.11, literal lines
9724--9864, including its transfinite construction and immediate proofs.
For the bounded multiplier and products ROOT freshly read the full
statements and immediate proofs in Sections 8.5--8.7, literal lines
10249--10490, and all of 8.8.1--8.8.7, lines 10494--10574.
Only the appropriate multiplier, positive-Laplacian and mixed-product
statements are used. In particular, no inference here uses the
retrieved literal-text Corollary 8.6.7 that an arbitrary positive continuous
multiplier preserves being a closed Bedford--Taylor current.
That statement cannot supply the closedness needed above; its PDF
pixels have not been freshly authenticated here.
ROOT also read the complete smooth Stokes/Green statement and
immediate proof, 5.11.1, lines 6990--7024. The deeper paralin\'eaire
Stokes proof 4.6.15 and all foundational positivity proofs are disclosed
background, not newly claimed full recursive proof-tree readings.
The proof above also uses directly the ordinary definition of a
closed current for integration of derivatives of smooth test forms.
For the maximum in bounded comparison ROOT freshly read the entire
statements and immediate proofs of 8.2.1--8.2.4, lines 9885--9980;
the soft-maximum lemma 1.7.6 cited there remains disclosed background.

The bounded branch calculus, raw mass, and two-atom calibration are
the unreviewed full drafts 383--387, with their own exact source
boundaries. ROOT's new complete local tropical dimension immediate
dependency reading and both full errata are separately recorded in
\url{P16.7-Ducros2012-local-dimension-ROOT-primary-reading-and-version-scope-20261009.json}.
That record documents actual author-version printed typographical
errors and version distinctions; it does not claim every Temkin or
ACVF foundational proof newly read. Full primary text, source pixels,
and book contexts remain LOCALONLY. No supplied or native program
was executed or imported.

The full unbounded intrinsic class, all-direction decreasing product
continuity and comparison, an appropriate separated intrinsic
compactness theorem, every obstacle/envelope and orthogonality step,
an arbitrary-support measure criterion, and full-class unbounded
existence and uniqueness have not been established by this paper.
Its discrete branch criterion is a scoped advance conditional on the
explicit geometric calibration, not an original-problem closure.
Compilation and independent review are pending. The protected P16.4
editor stays open and unchanged.
\end{document}
